\section{Discussion}
\label{sec:discussion}

\subsection{A general governance-generative substrate}
\label{sec:disc-substrate}

Across ten development milestones, the model demonstrated generative
sufficiency at every layer it was asked to support: bilateral cooperation
and reciprocity produce measurable, heterogeneous cooperation dynamics; an
accounted ecology produces genuine regime diversity; locally shared rules,
proposed and adopted from local information alone, reproducibly form and
persist unforced; delegation, discretion, and accountability all activate,
interact, and are each independently detectable in the observer layer;
institutions can observe each other, transmit components, compete via
migration, and scale across jurisdictions; and, as of Level 1.3/1.3b,
historically structured pressure, entrenchment, and legitimacy accumulate
and causally shape institutional response, sham-validated against a
fork-machinery artifact. None of this required, at any point, encoding a
target macro-state, a governance label, or a trajectory count. This is the
sense in which we call the substrate \textit{generatively sufficient}: it
is capable of producing institutions, institutional variation, and now
historically conditioned institutional trajectories, as genuine emergent
consequences of micro-rules.

\subsection{Institutional emergence without discrete phenotype convergence}
\label{sec:disc-emergence-without-convergence}

What the discovery pipeline found, at every scale and in both the
synchronic and historical substrates, is that generative sufficiency does
not, by itself, imply \textit{phenotypic convergence} onto a small number
of recurring configurations. Level 1.2's null result was legitimately
ambiguous, because near-universal collapse could have been suppressing
institutional variation the discovery pipeline never got to see. Level
1.2b removed that ambiguity and the result strengthened rather than
changed in kind: a \texttt{continuum} verdict, independently replicated as
\texttt{null\_replicated} on fresh data. Level 1.3 then asked whether
adding history changes this answer, on pooled, exploratory evidence
leaning \texttt{continuum} again but resting on one post hoc pool. Level
1.3b asked the same question with a purpose-designed, preregistered,
causally instrumented program, and the answer sharpened rather than
softened: \texttt{D\_persistent\_continuum}, with a discreteness rule that
failed to clear even its own entry-level silhouette floor at any tested
$k$ --- a more decisive rejection of discrete structure than the
synchronic substrate itself produced. History is real and causally
demonstrated; discreteness is not.

\subsection{History is real, causal, and persistent --- without producing discrete families}
\label{sec:disc-history-real}

This is the paper's own central, and initially counter-intuitive,
finding, and it is worth stating plainly because it is easy to conflate
the two questions it separates. ``Does history matter?'' and ``does
history sort outcomes into discrete families?'' are different questions,
and Level 1.3b's evidence answers them differently. History matters:
Section~\ref{sec:level13b-h5h6}'s sham-validated intervention design ---
the strongest causal design this project has built --- shows that erasing,
reversing, or exposure-equalizing a run's own realized pressure history
produces a real, non-decaying, and in several tests \textit{growing}
divergence from that same run with its history intact, validated against a
control proven to introduce zero spurious divergence on its own. This is
not a correlational association between history-adjacent variables and
outcomes; it is a controlled comparison of a system against a version of
itself with a specific piece of its own past removed. History does not,
however, sort outcomes into discrete families: the same milestone's own
discreteness test, applied via the identical unmodified rule Level 1.2b
established for the synchronic substrate, rejects discreteness more
decisively than the synchronic case did. A system can be genuinely,
causally path-dependent --- shaped by, and persistently different because
of, its own specific history --- while the space of trajectories that
history can produce remains continuous rather than partitioned into a
small number of recurring types. Increasing-returns and lock-in dynamics
in the historical-institutionalist literature
(\cref{sec:historical-institutionalism}) are typically illustrated with
discrete case comparisons (this country locked in one path, that country
another), but nothing in that theoretical tradition actually requires the
space of possible paths to be discrete; a continuous space with strong,
persistent path dependence is equally consistent with it, and is what this
model's own evidence supports.

\subsection{The richer generative state and the triangulation logic}
\label{sec:disc-triangulation}
\label{sec:disc-richer-state}

Level 1.2/1.2b's own observer-only phenotype vector characterized a
synchronic configuration. Level 1.3/1.3b's own historical mechanisms make
explicit that the object worth tracking is richer: a governance
configuration together with the history that produced it. We formalize
this, at the level of the model's own generative state, as

\begin{equation}
\mathcal{G}_t = [S_t,\, Z_t,\, H_t,\, E_t,\, L_t,\, P_t,\, R_t,\, C_t,\, \ldots]
\label{eq:rich-state}
\end{equation}

where $S_t$ and $Z_t$ denote the model's own observer-side state-capacity-
and constraint-adjacent composites ($S_{\text{obs}}$, $Z_{\text{obs}}$,
\cref{sec:observer-metrics-l13}) at simulation time $t$; $H_t$ is
accumulated historical/pressure state (the specific realized sequence of
security, bargaining, and coordination-complexity pressure, not merely its
current level); $E_t$ is accumulated entrenchment; $L_t$ is the dual
performance/procedural legitimacy state; $P_t$ denotes the currently active
pressure levels themselves (as distinct from their accumulated history
$H_t$); $R_t$ denotes realized institutional structure (delegation,
monitoring, sanctioning configuration); and $C_t$ denotes coordination/
transmission state across jurisdictions. The ellipsis is deliberate: this
is not a claim that eight dimensions are exhaustive, only that $S_t$ and
$Z_t$ alone --- the dissertation's own two axes --- are not sufficient to
characterize a governance configuration once history is admitted as
causally consequential. $S$ and $Z$ remain the central, interpretively
load-bearing axes of this expanded state: they are what the dissertation's
own SEM analysis identifies from real survey data, and what this paper's
own $S_{\text{obs}}$/$Z_{\text{obs}}$ composites are built, independently,
to approximate. The other components of $\mathcal{G}_t$ do not replace $S$
and $Z$; they explain \textit{movement through} $S$/$Z$ space --- why a
given configuration sits where it does and how it is likely to move next
--- in a way a purely synchronic, two-dimensional snapshot cannot.
\Cref{tab:expanded-state} summarizes every symbol in $\mathcal{G}_t$
against its concrete model implementation.

\begin{table}[htbp]
\centering
\caption{The expanded generative state $\mathcal{G}_t$: symbol, denotation, and model implementation.}
\label{tab:expanded-state}
\adjustbox{max width=\textwidth}{\begin{tabular}{@{}lp{6.5cm}p{5.8cm}@{}}
\toprule
\textbf{Symbol} & \textbf{Denotes} & \textbf{Model implementation} \\
\midrule
$S_t$ & State-capacity-adjacent observer composite & $S_{\text{obs}}$: mean of delegation, administrative\_function, monitoring, interpatch\_coordination, performance\_legitimacy \\
$Z_t$ & Institutional-constraint-adjacent observer composite & $Z_{\text{obs}}$: mean of accountability, bargaining\_participation, procedural\_legitimacy \\
$H_t$ & Accumulated historical/pressure state (full realized sequence, not just current level) & Per-channel pressure history series (security, bargaining, coordination-complexity) \\
$E_t$ & Accumulated institutional entrenchment & Age-dependent amendment/dissolution-resistance index \\
$L_t$ & Dual legitimacy state & Independent performance ($L_P$) and procedural ($L_R$) ledgers \\
$P_t$ & Currently active pressure levels (as distinct from accumulated history $H_t$) & Instantaneous per-channel pressure value at step $t$ \\
$R_t$ & Realized institutional structure & Delegation/monitoring/sanctioning configuration \\
$C_t$ & Coordination/transmission state across jurisdictions & Interpatch coordination links, component transmission history \\
$\ldots$ & Deliberately open --- not claimed exhaustive & \\
\bottomrule
\end{tabular}
}
\end{table}

This motivates the paper's own triangulation logic, introduced in
\cref{sec:intro-triangulation} and developed fully here. Three lines of
evidence bear on the dissertation's own $G=f(S,Z)$ framework, each
answering a question the others structurally cannot. \textit{SEM
identification} (the dissertation's own contribution) establishes, from
real cross-sectional survey data, that a two-factor $S$/$Z$ space fits
better than a collapsed one-dimensional ranking, and that real countries'
measured positions are not uniformly distributed within it. \textit{ABM
generative sufficiency} (Level 1.2/1.2b, extended by Level 1.3/1.3b)
establishes that a microscopic mechanism set, encoding none of $S$, $Z$, or
any trajectory label, is sufficient to produce heterogeneous configurations
occupying an analogous space, and --- new to this paper --- that adding
historically structured pressure and sequencing to that mechanism set
produces further, causally real heterogeneity without collapsing to a
small number of forms. \textit{Historical causal experimentation} (Level
1.3b's own sham-validated intervention design, unique to the ABM and
structurally unavailable to cross-sectional survey analysis, since no
survey can rerun a country's own history with one piece of it removed)
establishes that history's effect on institutional trajectory is causal,
not merely associational, within the tested mechanism set. None of these
three lines of evidence duplicates the others, and none is sufficient
alone: SEM identification cannot speak to generative sufficiency or
causal mechanism; ABM generative sufficiency cannot speak to whether real
countries' own historical trajectories are in fact governed by the same
dynamics; and historical causal experimentation within a simulated model
cannot, by itself, establish that real-world governance history behaves
the same way. Read together, they support a specific, bounded claim,
stated in full in \cref{sec:disc-boundary-claims} below. \Cref{tab:method-triangulation}
lays out all four sources of evidence side by side, including the
dissertation's own qualitative interpretive step, against what each does
and does not establish.

\begin{table}[htbp]
\centering
\caption{Methodological triangulation: what each line of evidence
establishes, its design, and its own limits.}
\label{tab:method-triangulation}
\adjustbox{max width=\textwidth}{\begin{tabular}{@{}p{2.6cm}p{4.0cm}p{4.0cm}p{4.2cm}@{}}
\toprule
\textbf{Method} & \textbf{What it establishes} & \textbf{Design} & \textbf{What it cannot establish} \\
\midrule
SEM identification (dissertation) & A 2-factor $S$/$Z$ space fits real survey data better than a 1-D ranking; real countries occupy non-uniform positions in it & Cross-sectional SEM, WVS Wave 7, real national cases & Generative sufficiency; causal mechanism; historical trajectory; discreteness of any typology laid over the space \\
ABM generative sufficiency (Level 1.2/1.2b, extended 1.3/1.3b) & A microscopic mechanism set, encoding none of $S$/$Z$/labels, is sufficient to produce heterogeneous configurations in an analogous space & Outcome-blind discovery, preregistered, independently replicated on fresh data & Whether real countries' own history follows the same dynamics; validation of the dissertation's specific empirical positions \\
Historical causal experimentation (Level 1.3b) & History's effect on simulated institutional trajectory is causal, not merely associational & Sham-controlled matched-seed intervention (state-copy control proven zero-divergence) & Whether real-world governance history behaves the same way; external validity beyond the tested mechanism set \\
Qualitative interpretation (dissertation \S4.4, this paper \S\ref{sec:disc-relationship-to-dissertation}) & A coherent narrative linking the two-dimensional space to named historical/institutional archetypes & Theorized typology over an identified coordinate space, not a cluster analysis & A statistical claim of discreteness; the dissertation's own text is explicit on this point (\cref{sec:intro-synchronic-to-historical}) \\
\bottomrule
\end{tabular}
}
\end{table}

\subsection{Continuum versus typology, restated for the historical substrate}
\label{sec:disc-continuum-vs-typology}

The methodological point from \cref{sec:why-null-informative} applies with
undiminished force to Level 1.3b's own eight-way result: silhouette-based
model selection alone cannot distinguish a genuine discrete structure from
a single dominant bisection of a smooth continuum, and a Bayesian
information criterion that is still falling at the edge of a tested range
is a signature of a still-improving flexible density fit, not evidence of
a converged discrete-family count. Level 1.3b's own BIC curve
(\cref{fig:l13b-bic}) falls sharply and without turning upward all the way
to $k{=}8$, the edge of the tested range --- precisely the pattern that
should raise, not settle, the question of whether a wider search would
find an even lower BIC at $k{=}9$ or beyond, which is exactly why this
paper does not report BIC in isolation anywhere and why
\cref{fig:l13b-bic} is captioned, and paired, with the weak cross-algorithm
and cross-stratum agreement evidence that argues directly against reading
$k{=}8$ as eight genuine institutional species (\cref{sec:disc-boundary-claims}).

\subsection{Relationship to the dissertation's S--Z trajectories}
\label{sec:disc-relationship-to-dissertation}

We can now state this relationship more precisely than Level 1.2b's own
discussion could, because Level 1.3b adds evidence the earlier milestone
did not have. The dissertation's SEM result is a statement about the joint
distribution of measured confidence and constraint variables across a
fixed set of existing national cases, at one historical moment, aggregated
at the level of already-formed states; its own text is explicit that this
identifies a coordinate space and does not classify countries or predict
trajectories, and that the three named trajectories are a qualitative,
theorized typology laid over that space, not a statistically discovered
partition (\cref{sec:intro-synchronic-to-historical}). This paper's own
ABM results, across all four milestones, are statements about what a
general microscopic mechanism set --- with and without historically
structured pressure --- can generate from scratch, blind to any expected
answer. Four specific things follow from combining them, stated precisely
to track exactly what was and was not tested.

First, the ABM's own generative substrate is sufficient to produce
heterogeneous governance configurations occupying a space with axes
plausibly analogous to $S$ and $Z$ (Levels 1.2/1.2b), and the observer-side
composites built for those axes do show a real, if modest and
weighting-sensitive, positive coupling (Level 1.3b Stage 9,
$r\approx0.26$--$0.29$ across defensible variants) --- weak, indirect
corroboration that state-capacity-like and constraint-like dimensions of a
generated governance space are related rather than orthogonal, consistent
with, though not a statistical test of, the dissertation's own finding
that $S$ and $Z$ are correlated but distinct factors. Second, historically
structured pressure and sequencing causally and persistently shape where a
simulated jurisdiction sits and how it got there (Level 1.3b Stage 6) ---
direct generative support for treating governance trajectories as
path-dependent processes, not merely static positions, which is exactly
the temporal dimension the dissertation's own SEM analysis is explicit
that it cannot supply (\cref{sec:historical-institutionalism}). Third,
neither the synchronic nor the historically extended substrate compresses
its own outcome space into three, or any other small fixed number of,
discrete, well-separated forms --- a null result about \textit{this
model's own generated space}, not a claim about whether real countries'
governance configurations are or are not discretely clustered, which this
paper's own ABM evidence cannot settle either way (\cref{sec:limitations}).
Fourth, and precisely because of the third point, the dissertation's own
three trajectories remain a legitimate, unaffected empirical finding at
their own analytical level: they are qualitative regions the dissertation
itself never claimed were statistically discrete, and this paper's own
result --- that a generative model producing an analogous space finds that
space continuous rather than discretely partitioned, even once
historically structured path dependence is added --- is consistent with,
and arguably strengthens the interpretive case for, reading the three
trajectories as coarse-grained, historically and institutionally
interpreted \textit{regions} of a continuous space rather than as discrete
natural kinds, which is a reading the dissertation's own methodology
(a two-dimensional coordinate space plus a qualitative typology, not a
cluster analysis) is already consistent with.

\subsection{What this paper does not claim: boundary conditions}
\label{sec:disc-boundary-claims}

Precisely because the evidence above is easy to overstate, we state five
specific boundary conditions explicitly, each ruling out a claim this
paper's own results could be mistaken for but do not support.

\textit{A statistical component count is not a species count.} Level
1.3b's leading discrete candidate has the lowest BIC at $k{=}8$ components;
this is a property of a specific model-selection criterion applied to a
specific finite sample, not a discovery of eight institutional species.
No tested $k$ clears even the classifier's own entry-level silhouette
floor (\cref{sec:level13b-h7h8}).

\textit{Statistical components are not recurrent institutional phenotypes.}
A recurrent phenotype, in this project's own vocabulary
(\cref{sec:terminology-discipline}), additionally requires passing
preregistered cross-seed, cross-algorithm, and bootstrap-stability
criteria; Level 1.3b's own $k{=}8$ candidate fails cross-algorithm
agreement (0.016--0.357 once the non-independent pair is excluded) and
cross-ecological-stratum reproduction ($-0.043$ to $-0.014$).

\textit{Recurrent phenotypes, had they been found, would not automatically
be dynamical attractors.} No claim anywhere in this paper asserts that any
discovered structure, at any milestone, is a state the system returns to
after perturbation; that specific, stronger claim requires direct evidence
of majority recovery on every tracked dimension
(\cref{sec:level12b-scale-causal}), which this project has never found for
any candidate structure at any milestone.

\textit{A multi-component model beating a single-Gaussian null is not
evidence of eight, or any specific number of, discrete families.} Level
1.3b's own Stage 8 fresh-data comparison shows the leading $k{=}8$
candidate beats a $k{=}1$ null on fresh data (\cref{sec:level13b-h7h8}) ---
but the leading \textit{continuous} candidate beats the same null by a
comparable margin, and an information criterion that has not converged
within the tested range cannot, on its own, license treating $k{=}8$ as
the true count rather than an artifact of where the search stopped.

\textit{Weak cross-algorithm and cross-stratum agreement rules out
treating any current partition as a discrete-family finding.} This is not
a subtle qualification; it is the single most decisive piece of evidence
in Level 1.3b's own record, and it is why every figure and table in this
paper reporting the $k{=}8$ result is paired with the cross-algorithm and
cross-stratum evidence that contextualizes it
(\cref{fig:l13b-bic,fig:l13b-discreteness}), never presented alone.

We add, finally, a fifth and more general boundary condition specific to
this paper's own claim structure: nowhere in this paper does a simulation
result from the ABM constitute empirical validation of the dissertation's
own three trajectories, of the exact number of governance archetypes real
countries occupy, or of the specific $S$/$Z$ values at which real countries
sit. The ABM's own $S_{\text{obs}}$/$Z_{\text{obs}}$ composites are
independently constructed, observer-side quantities, never fit to or
validated against the dissertation's own survey-derived latent factors;
any correspondence between them is, at most, weak structural corroboration
of the dissertation's own two-axis framing, not a validation of its
specific empirical content.

\subsection{Historical sequencing: confirmed as real, not merely motivated}
\label{sec:disc-historical-sequencing}

Level 1.2b's own discussion motivated, but explicitly did not test, a
formal decomposition $M_{\text{historical}} = M_{\text{general}} + H$. This
paper's second half tests it directly, and $H$ is no longer merely a
hypothesis: Level 1.3/1.3b implement $H$ as three pressure-transmission
channels, institutional entrenchment, and dual legitimacy ledgers, and
Level 1.3b's own sham-validated causal intervention design shows that $H$,
so implemented, causally and persistently shapes institutional response.
We restate \cref{eq:historical-decomposition-intro} with this paper's own
evidentiary status made explicit:

\begin{equation}
M_{\text{historical}} = M_{\text{general}} + H, \qquad
G_t = f\bigl(S_t, Z_t;\; H_{0:t},\, E_t,\, L_t,\, P_t,\, \ldots\bigr)
\label{eq:historical-decomposition-tested}
\end{equation}

where the macro-level right-hand side makes explicit that governance
outcome at time $t$ depends not only on current-state $S_t$, $Z_t$ but on
the full realized history $H_{0:t}$, accumulated entrenchment $E_t$,
legitimacy state $L_t$, and current pressure $P_t$ --- exactly the
dependency structure $\mathcal{G}_t$ (\cref{eq:rich-state}) is built to
track. What remains genuinely untested, and what this paper does not
claim, is whether $H$ as implemented here corresponds to the specific
historical mechanisms (critical junctures, colonial extraction, war-driven
state-building) the historical-institutionalist literature identifies in
real cases, whether the model's own parameter domain and duration are
representative of real historical timescales, and whether real governance
systems' own historical processes are governed by dynamics of the same
qualitative shape this model's tested mechanism set produces. Explaining
why real governance systems concentrate around particular recurring
trajectories --- as opposed to why a generative model of the kind built
here does or does not --- is deferred to future, historically calibrated
research (\cref{sec:limitations}).
